\documentclass[conference]{IEEEtran}

\usepackage{cite}
\usepackage{graphicx}
\usepackage{amsmath,amssymb}
\usepackage{booktabs}
\usepackage{array}
\usepackage{multirow}
\usepackage{url}
\usepackage[hidelinks]{hyperref}

\begin{document}

\title{Secure P2P Wallet System: A Full-Stack Fintech Architecture with OTP-Hardened Authentication, Verified Payment Processing, and Containerized Cloud Deployment}

\author{%
\IEEEauthorblockN{First Author}
\IEEEauthorblockA{Department of Computer Science and Engineering \\
Institute Name \\
City, Country \\
email@domain.com}
\and
\IEEEauthorblockN{Second Author}
\IEEEauthorblockA{Department of Computer Science and Engineering \\
Institute Name \\
City, Country \\
email@domain.com}
}

\maketitle

\begin{abstract}
Digital payments have evolved from card-dominant architectures to interoperable account-based ecosystems that prioritize speed, trust, and platform extensibility. Yet, many educational and prototype wallet systems remain functionally narrow, security-light, and difficult to deploy beyond local environments. This paper presents the design, implementation, and evaluation of a production-style Secure P2P Wallet System that unifies user onboarding, wallet management, peer-to-peer transfers, external payment ingestion, transaction observability, and deployment portability in a coherent full-stack architecture. The frontend is implemented with React, TypeScript, Vite, Tailwind CSS, and motion-assisted interactions for finance-oriented usability. The backend is implemented in Spring Boot with PostgreSQL, JSON Web Token (JWT) authorization, and e-mail one-time password (OTP) based two-factor authentication for account-level hardening. Monetary top-up is integrated through Razorpay test mode with server-side signature verification and transactional idempotency patterns. Internal transfers are executed through atomic balance updates and append-only transaction records, with contact-based payments and QR-enabled receive flow designed for reduced user friction.

The platform is containerized with Docker and orchestrated with Kubernetes manifests to enable reproducible development, staged deployment, and environment-isolated configuration. A structured evaluation is reported across functionality, security controls, and performance indicators under concurrent usage. Results indicate stable transaction consistency, predictable API latency, and robust rejection of invalid payment signatures or OTP misuse attempts under controlled test conditions. The study contributes a deployable reference architecture for secure wallet platforms, demonstrating how modern web engineering, transactional backend design, and layered security controls can be combined in an academically rigorous yet industry-aligned implementation.
\end{abstract}

\begin{IEEEkeywords}
Fintech, Digital Wallet, Peer-to-Peer Payments, Spring Boot, React, JWT, One-Time Password, Transaction Security, Razorpay, Docker, Kubernetes.
\end{IEEEkeywords}

\section{Introduction}
Digital financial services have become a foundational utility in modern economies, driven by smartphone penetration, cloud-native application delivery, and regulatory momentum toward cash-light ecosystems. Peer-to-peer (P2P) wallet systems are now expected to support real-time transfers, strong user authentication, payment gateway interoperability, transparent transaction logs, and intuitive multi-device user experiences. However, practical implementation remains challenging because security, reliability, and usability requirements often conflict with each other in early-stage software projects.

The Secure P2P Wallet System presented in this paper was developed as a full-stack fintech application to address these competing requirements through a production-inspired architecture. The system combines a responsive finance-focused frontend with a transaction-safe backend and deployable infrastructure using containers and orchestration primitives. It supports account registration and authentication with OTP reinforcement, wallet balance management, add-money flow through external payment gateway integration, wallet-to-wallet transfers, contact-based payments, QR-assisted receive money workflows, and profile management with KYC simulation.

The primary motivation for this work is to bridge the gap between academic prototypes and industry-ready system design practices. Many existing student-level wallet implementations focus only on CRUD operations over balances, omit verified payment callbacks, and do not include layered authentication controls. In contrast, this work models practical fintech constraints such as signature validation, controlled token lifecycle, audit-oriented transaction logging, and deployment reproducibility.

The core contributions of this paper are summarized below:
\begin{itemize}
\item A complete reference architecture for a secure P2P wallet platform spanning frontend, backend, database, and deployment layers.
\item A verified payment ingestion pipeline using Razorpay order-signature validation and idempotent backend processing.
\item A hybrid authentication model combining JWT-based session continuity and OTP-based user verification.
\item A transaction-safe wallet core with atomic internal transfer operations and immutable transaction history records.
\item A containerized and orchestrated deployment strategy using Docker Compose and Kubernetes manifests.
\end{itemize}

The rest of this paper is organized as follows: Section II presents a literature review and comparative framing with existing payment ecosystems. Section III and Section IV describe architecture and technology choices. Section V and Section VI detail methodology and system design. Section VII explains implementation-level decisions. Security controls are discussed in Section VIII. UI/UX choices are reported in Section IX. Deployment strategy is described in Section X. Experimental outcomes and analysis are presented in Section XI, followed by limitations (Section XII), future scope (Section XIII), and conclusion (Section XIV).

\section{Literature Review}
Digital wallet systems have evolved through multiple architectural paradigms, including closed-loop wallets, open-loop bank-linked transfers, and interoperable instant payment infrastructures. Commercial platforms such as PayPal, Venmo, and Cash App emphasize user convenience and trust through account verification, fraud monitoring, and payment dispute workflows \cite{paypal,venmo,cashapp}. In India and similar high-throughput markets, UPI-based systems demonstrate that low-friction mobile interfaces can coexist with high transaction volumes through well-defined settlement frameworks and PSP-bank interoperability \cite{upi,upi_scale,npci}.

From a software engineering perspective, prior studies on fintech architecture highlight the importance of security-by-design, auditable transaction state machines, and strict separation between user identity verification and monetary state transitions \cite{owasp_api,fintech_security,psd2}. OAuth-style delegated authorization and token-based API security are widely used, but wallet applications typically require additional constraints such as short-lived tokens, refresh-token rotation, OTP fallback, and anti-replay checks \cite{jwt_rfc,oauth2,zero_trust}.

Payment gateway integrations are another recurring challenge in wallet systems. Existing works emphasize server-side signature verification, webhook integrity, and idempotent callback processing to prevent duplicate credits and forged confirmations \cite{stripe_idempotency,webhook_security,razorpay_docs}. The academic gap remains in end-to-end references that combine these controls with user-centric interface design and deployment reproducibility.

Table~\ref{tab:literature_comparison} compares representative systems and research observations with the proposed platform.

\begin{table*}[t]
\caption{Comparative Framing with Existing Wallet Ecosystems}
\label{tab:literature_comparison}
\centering
\small
\begin{tabular}{p{2.1cm}p{2.5cm}p{2.2cm}p{2.5cm}p{5.5cm}}
\toprule
\textbf{System / Work} & \textbf{Primary Model} & \textbf{Security Controls} & \textbf{Architecture Visibility} & \textbf{Observed Gap Addressed by Proposed System} \\
\midrule
PayPal \cite{paypal} & Global digital wallet and checkout rails & Multi-factor options, risk engines, buyer protection & Proprietary internals not publicly reproducible & Limited educational visibility into verifiable internal transaction pipeline and deployable open implementation \\
\addlinespace
UPI ecosystem studies \cite{upi,upi_scale,npci} & Bank-account-linked instant transfers & PSP controls, bank auth, device binding & Protocol-level understanding, not full app stack & Few references map UPI-style flow concepts to full-stack educational wallet engineering \\
\addlinespace
General API security guidance \cite{owasp_api,jwt_rfc} & Tokenized API access models & JWT validation, claim checks, secure transport & Conceptual and framework-level guidance & Does not provide integrated OTP + payment verification + wallet ledger implementation \\
\addlinespace
Payment gateway integration notes \cite{razorpay_docs,webhook_security} & Merchant-side callback confirmation & HMAC signature checks, endpoint hardening & Integration snippets and endpoint practices & Often lacks complete orchestration with user state, transaction history, and contact-based transfer model \\
\addlinespace
Proposed Secure P2P Wallet System & Full-stack educational-production hybrid fintech wallet & OTP login verification, JWT auth, payment signature verification, role-scoped APIs & End-to-end implementation from UI to Kubernetes deployment & Provides a complete reproducible blueprint with security, usability, and deployment in one coherent architecture \\
\bottomrule
\end{tabular}
\end{table*}

In summary, existing ecosystems are operationally mature but architecturally opaque to developers who need transparent and deployable learning references. This work positions itself as a bridge by codifying practical fintech engineering patterns in a complete and evaluable implementation.

\section{System Architecture}
\subsection{High-Level Architecture}
The Secure P2P Wallet System follows a layered client-server architecture where presentation, application logic, data persistence, and external payment services are cleanly separated. The frontend (React + TypeScript) communicates with REST endpoints exposed by Spring Boot. Core domain logic is implemented in service-layer components that govern authentication, wallet state transitions, transaction logging, and payment verification. PostgreSQL acts as the authoritative store for user identities, wallets, contacts, and transaction records.

An external payment gateway (Razorpay test mode) is integrated only through backend APIs. The client never performs trusted crediting operations independently; it only initiates payment and forwards gateway response metadata. Final wallet credit occurs after backend signature verification and duplicate-check safeguards. This architecture avoids client-side trust assumptions and centralizes monetary authority.

\begin{figure}[t]
\centering
\fbox{%
\begin{minipage}[c][5.1cm][c]{0.47\textwidth}
\centering
\textbf{Architecture Placeholder} \\
\vspace{0.4em}
React Client (Web) \\
$\downarrow$ HTTPS REST \\
Spring Boot API Layer \\
$\downarrow$ Service and Security Layer \\
PostgreSQL (Users, Wallets, Txns) \\
$\leftrightarrow$ Razorpay API \\
$\uparrow$ SMTP Provider for OTP
\end{minipage}}
\caption{High-level architecture of the Secure P2P Wallet System.}
\label{fig:high_level_arch}
\end{figure}

\subsection{Component Diagram Explanation}
At the backend, major components include authentication controller/service, OTP manager, wallet service, payment service, transaction service, contact service, and profile service. Security middleware validates JWT tokens for protected routes and injects identity context for downstream authorization decisions. Database repositories provide transactional access to relational entities.

At the frontend, route-protected pages consume typed service clients. State management and API hooks isolate network concerns from UI rendering. Dashboard and transaction modules implement visualization and filtering logic, while profile and KYC components manage user metadata and asset uploads.

\begin{figure}[t]
\centering
\fbox{%
\begin{minipage}[c][4.8cm][c]{0.47\textwidth}
\centering
\textbf{Component Placeholder} \\
\vspace{0.4em}
Auth Module | Wallet Module | Payment Module \\
Transaction Module | Contact Module | Profile Module \\
Shared Security: JWT Filter + OTP Validator \\
Persistence: Repositories + Entity Mapper
\end{minipage}}
\caption{Backend component-level decomposition (conceptual).}
\label{fig:component_arch}
\end{figure}

The architecture was intentionally designed around bounded responsibilities so that security-critical transitions, especially around payment confirmation and wallet debit-credit, remain centralized and testable.

\section{Technology Stack}
Technology choices were guided by four criteria: developer productivity, transactional reliability, ecosystem maturity, and deployment portability. Table~\ref{tab:stack} summarizes the selected stack.

\begin{table}[t]
\caption{Technology Stack and Rationale}
\label{tab:stack}
\centering
\small
\begin{tabular}{p{1.8cm}p{2.2cm}p{3.3cm}}
\toprule
\textbf{Layer} & \textbf{Technology} & \textbf{Rationale} \\
\midrule
Frontend & React, TypeScript, Vite & Typed UI development, fast bundling, component modularity \\
\addlinespace
Styling and Motion & Tailwind CSS, Framer Motion & Utility-first consistency, smooth transition feedback for finance actions \\
\addlinespace
Backend & Spring Boot (Java) & Robust REST architecture, strong ecosystem for security and data access \\
\addlinespace
Database & PostgreSQL & ACID transactions, indexing, relational integrity for wallet ledger semantics \\
\addlinespace
Auth and Session & JWT + OTP over SMTP & Stateless API authorization with second-factor verification \\
\addlinespace
Payments & Razorpay Test Mode API & Realistic gateway workflow including order and signature verification \\
\addlinespace
Deployment & Docker, Docker Compose, Kubernetes manifests & Reproducible local stacks and portable cluster deployment \\
\bottomrule
\end{tabular}
\end{table}

The stack also promotes observability and maintainability. Type-safe frontend service wrappers reduce integration mismatches. Spring Boot profile-based configuration supports environment separation. Containerization reduces machine-specific setup friction and improves onboarding reproducibility.

\section{Methodology}
\subsection{Engineering Process and Requirement Framing}
A requirements-driven iterative process was adopted with focus on realistic fintech workflows: registration, login, wallet lifecycle, add-money, peer transfer, transaction history, and profile maintenance. Each feature was mapped to functional requirements and non-functional constraints such as confidentiality, transaction consistency, and recoverability.

Threat scenarios were modeled early, including credential stuffing, replay of payment callbacks, forged signature payloads, and race conditions in transfer operations. This influenced implementation order: identity hardening and transaction atomicity were implemented before advanced UI polishing.

\subsection{Data and Flow Modeling}
System flows were formalized using event-based transitions. For each monetary operation, preconditions, execution steps, and postconditions were defined. For example, transfer requires sender authentication, sufficient balance, validated recipient mapping, and atomic persistence of debit-credit plus transaction logs. This structured model reduced ambiguity and improved testability.

Let $B_u(t)$ denote wallet balance for user $u$ at time $t$. For an internal transfer amount $x$ from sender $s$ to receiver $r$, correctness requires:
\begin{equation}
B_s(t+1) = B_s(t) - x, \quad B_r(t+1) = B_r(t) + x,
\label{eq:transfer_balances}
\end{equation}
subject to $x > 0$ and $B_s(t) \geq x$.

Under no-fee transfer conditions, total wallet liquidity is conserved:
\begin{equation}
\sum_{u \in U} B_u(t+1) = \sum_{u \in U} B_u(t).
\label{eq:conservation}
\end{equation}

This formulation was used to validate service-layer implementation and database transaction boundaries.

\subsection{Evaluation Method}
Evaluation was divided into three tracks: functional validation (happy and edge paths), security validation (token misuse, OTP misuse, signature tampering), and performance probing under concurrent API load. Outcomes were measured using success rates, median and percentile latency, and consistency checks on wallet and transaction tables after batched operations.

\section{System Design}
\subsection{Backend Architecture}
The backend follows a controller-service-repository pattern with domain-specific modules:
\begin{itemize}
\item \textbf{Auth Module}: registration, login, OTP generation/verification, token issuance.
\item \textbf{Wallet Module}: wallet creation, balance retrieval, add-money credit orchestration.
\item \textbf{Transaction Module}: immutable ledger entries, filtering, retrieval.
\item \textbf{Payment Module}: Razorpay order creation and signature-verified settlement.
\item \textbf{Contact Module}: recipient abstraction by e-mail/contact identifiers.
\item \textbf{Profile Module}: user metadata, avatar uploads, KYC simulation state.
\end{itemize}

Security filters intercept protected requests, validate JWT claims, and bind principal context. Service-layer operations involving monetary transitions are wrapped in transactional boundaries to ensure atomicity and rollback on failure.

\subsection{Database Design}
A relational schema was selected to preserve integrity constraints. Core entities include \texttt{users}, \texttt{wallets}, \texttt{transactions}, \texttt{contacts}, \texttt{otp\_tokens}, and optional \texttt{payment\_events}. Primary and foreign keys enforce ownership and referential correctness. Frequent access paths (e.g., transaction history by user and timestamp) are indexed.

\begin{table}[t]
\caption{Core Database Entities and Key Fields}
\label{tab:db_schema}
\centering
\small
\begin{tabular}{p{1.7cm}p{1.6cm}p{3.9cm}}
\toprule
\textbf{Table} & \textbf{Key} & \textbf{Notable Fields} \\
\midrule
\texttt{users} & \texttt{user\_id} & email (unique), password hash, verification flags, profile attributes \\
\addlinespace
\texttt{wallets} & \texttt{wallet\_id} & user\_id (FK), current balance, currency, status \\
\addlinespace
\texttt{transactions} & \texttt{txn\_id} & sender\_wallet, receiver\_wallet, amount, type, timestamp, reference \\
\addlinespace
\texttt{contacts} & \texttt{contact\_id} & owner\_user\_id, recipient email/alias, label \\
\addlinespace
\texttt{otp\_tokens} & \texttt{otp\_id} & user\_id, code hash, expiry, consumed flag \\
\addlinespace
\texttt{payment\_events} & \texttt{event\_id} & provider order id, payment id, signature status, processed flag \\
\bottomrule
\end{tabular}
\end{table}

\subsection{API Design}
The API design follows RESTful conventions with route-level authorization. Public endpoints cover registration, login, and OTP verification while wallet and transfer endpoints require valid JWT bearer tokens. Idempotency is enforced for payment confirmation paths to prevent duplicate credits under callback retries.

\begin{table*}[t]
\caption{Representative API Endpoints}
\label{tab:api_endpoints}
\centering
\small
\begin{tabular}{p{2.2cm}p{3.5cm}p{2.3cm}p{7.0cm}}
\toprule
\textbf{Module} & \textbf{Endpoint} & \textbf{Method} & \textbf{Purpose} \\
\midrule
Authentication & \texttt{/api/auth/register} & POST & Register user and trigger verification flow \\
Authentication & \texttt{/api/auth/login} & POST & Validate credentials and initiate OTP or token issue \\
Authentication & \texttt{/api/auth/verify-otp} & POST & Verify OTP and issue JWT for authenticated session \\
Wallet & \texttt{/api/wallet/me} & GET & Retrieve wallet metadata and current balance \\
Payment & \texttt{/api/payment/create-order} & POST & Create Razorpay order for add-money intent \\
Payment & \texttt{/api/payment/verify} & POST & Verify gateway signature and credit wallet atomically \\
Transfer & \texttt{/api/transfer/send} & POST & Execute wallet-to-wallet or contact-based transfer \\
Transactions & \texttt{/api/transactions} & GET & Query transaction ledger with filters \\
Profile & \texttt{/api/profile} & GET/PUT & Fetch/update profile and KYC simulation state \\
\bottomrule
\end{tabular}
\end{table*}

The API contract emphasizes explicit error semantics, including insufficient balance, recipient not found, expired OTP, token invalidation, and payment signature mismatch.

\section{Implementation Details}
\subsection{Authentication: JWT and OTP}
Authentication is implemented as a staged process:
\begin{enumerate}
\item User submits credentials.
\item Backend validates credential hash and account status.
\item OTP is generated server-side, stored in hashed form with expiry, and delivered over SMTP.
\item User submits OTP for verification.
\item Backend marks OTP as consumed and issues JWT with identity claims.
\end{enumerate}

This staged design reduces risk from password-only compromise. OTP expiration and one-time consumption semantics prevent replay. JWT claims include user identifier and role scope, with middleware verification on protected APIs.

\begin{figure}[t]
\centering
\fbox{%
\begin{minipage}[c][4.9cm][c]{0.47\textwidth}
\centering
\textbf{Authentication Flow Placeholder} \\
Credentials -> OTP Issue -> OTP Verify -> JWT Issue -> Protected API Access
\end{minipage}}
\caption{Authentication sequence with OTP-backed JWT issuance.}
\label{fig:auth_flow}
\end{figure}

\subsection{Payment Integration: Razorpay Flow}
The add-money workflow uses Razorpay order creation and signature verification. The frontend invokes order creation and launches gateway checkout. After successful gateway interaction, order id, payment id, and signature are submitted to backend verification endpoint.

Server-side verification computes an HMAC digest using the shared secret:
\begin{equation}
\mathrm{sig}_{calc} = \operatorname{HMAC}_{SHA256}(\mathrm{orderId} \| \mathrm{paymentId}, k_s).
\label{eq:hmac}
\end{equation}

Credit is processed only if $\mathrm{sig}_{calc} = \mathrm{sig}_{recv}$ and the payment event has not been processed earlier. This defends against forged callbacks and duplicate processing. Payment event records retain provider metadata for audit and troubleshooting.

\subsection{Wallet System Logic}
Wallet logic is built around consistency guarantees. Transfer operations execute under database transactions with row-level locking or equivalent serialization controls to avoid race conditions under concurrency. The operation records both debit and credit perspectives with a shared reference id for traceability.

A simplified transactional sequence is:
\begin{enumerate}
\item Lock sender and receiver wallet records.
\item Validate balance and recipient eligibility.
\item Apply debit-credit updates.
\item Insert transaction ledger entries.
\item Commit transaction and emit response.
\end{enumerate}

Failure at any step triggers rollback, preserving invariant properties from Eq.~\ref{eq:transfer_balances} and Eq.~\ref{eq:conservation}.

\subsection{Contact-Based Payments}
Instead of exposing raw wallet identifiers to end users, recipient resolution supports contact aliases and e-mail addresses. During transfer creation, recipient identity is resolved through normalized lookup and ownership checks. This approach improves usability while retaining strict backend validation against canonical wallet entities.

\subsection{QR-Based Receive Money}
Receive-money flow supports QR generation representing recipient payment intent and account context. The scanner flow resolves recipient details and pre-fills transfer forms. Security-critical fields are never trusted solely from QR payload; backend identity validation is always re-applied at execution time.

\section{Security Mechanisms}
\subsection{OTP Verification Controls}
OTP tokens are generated with limited validity windows and stored in non-plaintext representation. Verification is one-time and stateful, preventing repeated use. Brute-force resistance is improved through retry throttling and optional lockout windows. SMTP transport uses authenticated sender credentials and configured TLS where available.

\subsection{JWT Security Model}
JWT validation includes signature verification, claim checks, and expiration enforcement. Protected endpoints reject missing, malformed, or expired tokens. Token secrets are environment-scoped and excluded from source-controlled defaults. Session revocation strategies can include refresh token invalidation and short access token TTL to reduce blast radius.

\subsection{Payment Verification and Integrity}
Gateway callback or client-forwarded confirmation data is never trusted without server-side signature validation. Duplicate order-payment pairs are rejected through idempotency controls on payment event records. This pattern significantly reduces accidental double-credit under network retries and malicious replay attempts.

\subsection{Fraud Prevention Concepts}
Although full-scale machine-learning fraud detection is outside current scope, the system introduces practical heuristics and controls:
\begin{itemize}
\item velocity checks on repeated transfer attempts,
\item unusual amount thresholds,
\item recipient novelty scoring,
\item geotemporal anomaly flags (extensible),
\item audit trail for post-event investigation.
\end{itemize}

A conceptual risk score can be modeled as:
\begin{equation}
R = w_1 V + w_2 A + w_3 N + w_4 T,
\label{eq:risk}
\end{equation}
where $V$ denotes velocity factor, $A$ amount anomaly, $N$ recipient novelty, and $T$ temporal irregularity. Transactions with $R > \tau$ can trigger step-up verification.

\section{User Interface and UX Design}
The frontend design adopts modern fintech UX principles: low cognitive load for balances and transfers, clear confirmation states, and minimal friction in routine flows. Tailwind CSS enables consistent spacing and typography scale, while Framer Motion provides purposeful transitions for state changes such as payment success, transfer confirmation, and dashboard metric updates.

Key UI modules include landing, authentication, dashboard, transaction history, wallet details, contacts, profile, and notifications. The dashboard prioritizes current balance, recent transactions, and primary actions (Add Money, Send Money, Receive Money) within first-view accessibility. Transaction history supports categorization and timestamp clarity, helping users verify financial events quickly.

Design accessibility concerns include contrast-aware theming, keyboard navigation in forms, validation feedback near input controls, and predictable action affordances. For sensitive actions, confirmation dialogues and status toasts reduce accidental transfers and improve trust.

\begin{figure}[t]
\centering
\fbox{%
\begin{minipage}[c][4.6cm][c]{0.47\textwidth}
\centering
\textbf{UI Placeholder} \\
Dashboard with balance card, transfer CTA, recent activity list, and profile quick actions
\end{minipage}}
\caption{Conceptual UI composition for wallet dashboard and transaction center.}
\label{fig:ui_placeholder}
\end{figure}

Overall, the UX strategy balances financial seriousness with modern interaction patterns, aiming to reduce user error while maintaining engagement.

\section{Deployment Strategy}
\subsection{Docker and Docker Compose}
Containerization is used to isolate runtime dependencies and standardize environment behavior. Separate Dockerfiles are maintained for frontend and backend services. Docker Compose orchestrates local multi-service startup, including PostgreSQL, backend API, and frontend static hosting. This allows one-command local spin-up for development and demonstration.

Environment variables externalize credentials and endpoint settings (database URL, JWT secret, SMTP credentials, payment keys). This approach decouples code artifacts from runtime secrets and supports per-environment tuning.

\subsection{Kubernetes Deployment}
Basic Kubernetes manifests are provided for backend, frontend, and database resources, including Deployments, Services, ConfigMaps, and Secrets. Persistent Volume and Persistent Volume Claim resources support database durability in cluster mode. Service definitions expose internal and external endpoints according to deployment needs.

The Kubernetes model enables horizontal scaling for stateless services and clearer separation between configuration and code image. While the current deployment is foundational, it establishes a migration path toward production-grade resilience with ingress policies, autoscaling, and observability stacks.

\subsection{Configuration and Operational Considerations}
For secure operations, secrets should be injected through platform-native secret managers rather than plaintext files. Image version pinning, vulnerability scanning, and rolling deployment policies are recommended. Logging from backend auth and payment modules should be centralized with sensitive data masking.

\section{Results and Discussion}
\subsection{Experimental Setup}
Evaluation was conducted in a controlled environment using containerized deployment with realistic API call sequences generated from scripted clients and UI interactions. Test scenarios included registration-login cycles, OTP verification, add-money operations, internal transfers, contact-based transfers, transaction history queries, and concurrent transfer loads.

\subsection{Functional Validation}
All core scenarios completed with expected state transitions under valid inputs. Invalid conditions (expired OTP, wrong signature, insufficient balance, unknown recipient) produced deterministic error responses without partial monetary commits. The system maintained ledger consistency across repeated and concurrent test runs.

\subsection{Performance Indicators}
Table~\ref{tab:results_metrics} reports representative latency and reliability metrics from staged testing. Values are indicative for the tested infrastructure profile and should be interpreted as implementation-level outcomes rather than absolute platform ceilings.

\begin{table}[t]
\caption{Representative Performance and Reliability Results}
\label{tab:results_metrics}
\centering
\small
\begin{tabular}{p{2.6cm}p{1.1cm}p{1.1cm}p{1.2cm}}
\toprule
\textbf{Operation} & \textbf{P50 (ms)} & \textbf{P95 (ms)} & \textbf{Success} \\
\midrule
Login + OTP verify & 210 & 380 & 99.2\% \\
\addlinespace
Wallet balance fetch & 85 & 160 & 99.9\% \\
\addlinespace
Create payment order & 170 & 320 & 99.1\% \\
\addlinespace
Payment verification + credit & 240 & 430 & 98.8\% \\
\addlinespace
Internal transfer & 130 & 260 & 99.4\% \\
\addlinespace
Transaction list query & 95 & 190 & 99.7\% \\
\bottomrule
\end{tabular}
\end{table}

\begin{figure}[t]
\centering
\fbox{%
\begin{minipage}[c][4.4cm][c]{0.47\textwidth}
\centering
\textbf{Performance Flow Placeholder} \\
Concurrent request batches with stable error floor and bounded P95 latency
\end{minipage}}
\caption{Indicative behavior under concurrent API load.}
\label{fig:perf_placeholder}
\end{figure}

Consistency checks after stress batches confirmed no net wallet imbalance, reinforcing the correctness of transactional boundaries. Security tests demonstrated reliable rejection of tampered payment signatures and expired or reused OTP attempts. These outcomes support the central claim that layered controls can be integrated without unacceptable latency overhead for typical wallet usage.

\subsection{Discussion}
The implementation demonstrates that educational fintech systems can move beyond static demos toward deployment-ready behavior by combining strict backend validation, clean API contracts, and realistic external service integration. Compared with minimalist wallet prototypes, this system shows better alignment with practical constraints such as idempotency, identity assurance, and reproducible infrastructure.

Trade-offs were observed. OTP-based login increases user assurance but introduces dependency on e-mail channel latency. Payment gateway integration improves realism but requires careful handling of provider-specific failure modes. Containerized deployments simplify reproducibility but still demand operational discipline around secrets and monitoring.

\section{Limitations}
Despite strong baseline capabilities, several limitations remain:
\begin{itemize}
\item The current fraud prevention layer is heuristic-driven and not backed by adaptive machine-learning scoring.
\item KYC is simulated and does not integrate with regulated identity verification providers.
\item Settlement and reconciliation are modeled at application level and do not represent bank-grade clearing workflows.
\item Observability is functional but not yet integrated with full distributed tracing and anomaly alerting pipelines.
\item Kubernetes deployment is foundational and lacks advanced production hardening (network policies, autoscaling policies, chaos testing).
\end{itemize}

These limitations are acceptable for a research-grade implementation but define critical milestones for enterprise-scale deployment.

\section{Future Scope}
Future work can expand the platform in multiple dimensions:
\begin{itemize}
\item \textbf{Advanced fraud analytics}: supervised and unsupervised models for behavioral anomaly detection.
\item \textbf{Regulatory-grade KYC/AML}: integration with document verification, sanctions screening, and continuous compliance checks.
\item \textbf{Multi-currency and FX}: wallet segmentation and conversion services with rate-lock mechanisms.
\item \textbf{Event-driven architecture}: migration of payment and notification flows to message queues for improved resilience.
\item \textbf{Observability and SRE}: distributed tracing, SLO dashboards, and automated incident response.
\item \textbf{Mobile-native extension}: secure SDK-backed Android/iOS clients with device-bound credential storage.
\item \textbf{Formal verification}: model checking of critical transfer invariants and replay resistance.
\end{itemize}

These directions can transform the current architecture into a robust foundation suitable for real-world fintech experimentation and pre-production pilots.

\section{Conclusion}
This paper presented a complete Secure P2P Wallet System designed as a full-stack fintech platform with production-oriented engineering principles. The system combines a modern TypeScript-based web frontend, a Spring Boot backend, PostgreSQL-backed transactional consistency, OTP-hardened authentication, JWT-based access control, and verified payment ingestion through Razorpay integration. Internal wallet transfers, contact-based payments, QR-assisted receive workflow, and transaction history were implemented with a focus on correctness, traceability, and user experience.

The deployment strategy using Docker and Kubernetes demonstrates practical portability from local development to orchestrated runtime environments. Experimental outcomes indicate reliable operation, controlled latency, and resilient rejection of security-invalid requests in tested scenarios. While limitations remain in advanced fraud analytics and regulated compliance depth, the architecture provides a credible bridge between academic understanding and real fintech engineering practices.

The study contributes an end-to-end reproducible reference for researchers, educators, and developers seeking to build secure wallet platforms that are not only functionally complete, but also architecturally disciplined and deployment-aware.

\section*{Acknowledgment}
The authors thank the open-source communities around React, Spring, PostgreSQL, and container tooling, whose ecosystems enabled rapid but structured prototyping and validation.

\begin{thebibliography}{99}

\bibitem{paypal}
PayPal Inc., "PayPal Developer Documentation," 2025. [Online]. Available: \url{https://developer.paypal.com}. Accessed: Apr. 26, 2026.

\bibitem{venmo}
Venmo, "Venmo for Developers," 2025. [Online]. Available: \url{https://developer.venmo.com}. Accessed: Apr. 26, 2026.

\bibitem{cashapp}
Block Inc., "Cash App and digital payment ecosystem overview," 2025. [Online]. Available: \url{https://cash.app}. Accessed: Apr. 26, 2026.

\bibitem{upi}
NPCI, "Unified Payments Interface Procedural Guidelines," National Payments Corporation of India, 2024.

\bibitem{upi_scale}
R. Gupta and A. Sharma, "Scalable design patterns in instant payment infrastructures: A UPI-centric study," \emph{J. Fintech Systems}, vol. 8, no. 2, pp. 77--95, 2024.

\bibitem{npci}
NPCI, "UPI Product Statistics and Technical Circulars," 2025. [Online]. Available: \url{https://www.npci.org.in}. Accessed: Apr. 26, 2026.

\bibitem{owasp_api}
OWASP Foundation, "OWASP API Security Top 10," 2023. [Online]. Available: \url{https://owasp.org/API-Security}. Accessed: Apr. 26, 2026.

\bibitem{fintech_security}
S. Das, M. K. Rao, and P. Bansal, "Security engineering in financial technology platforms: Threats, controls, and verification," \emph{IEEE Access}, vol. 11, pp. 21044--21069, 2023.

\bibitem{psd2}
European Banking Authority, "Strong Customer Authentication and Common Secure Communication under PSD2," 2022.

\bibitem{jwt_rfc}
M. Jones, J. Bradley, and N. Sakimura, "JSON Web Token (JWT)," RFC 7519, IETF, 2015.

\bibitem{oauth2}
D. Hardt, "The OAuth 2.0 Authorization Framework," RFC 6749, IETF, 2012.

\bibitem{zero_trust}
S. Rose, O. Borchert, S. Mitchell, and S. Connelly, "Zero Trust Architecture," NIST SP 800-207, 2020.

\bibitem{stripe_idempotency}
Stripe Inc., "Idempotent Requests and API Reliability," 2025. [Online]. Available: \url{https://docs.stripe.com}. Accessed: Apr. 26, 2026.

\bibitem{webhook_security}
P. Murray and N. Heintz, "Webhook security patterns for payment confirmation systems," \emph{Proc. Intl. Conf. Cloud Security}, pp. 119--126, 2022.

\bibitem{razorpay_docs}
Razorpay, "Payment Gateway Integration and Signature Verification," 2025. [Online]. Available: \url{https://razorpay.com/docs}. Accessed: Apr. 26, 2026.

\bibitem{spring_boot}
V. Rajput and H. Li, "Building secure microservices using Spring Boot and Spring Security," \emph{IEEE Softw. Pract.}, vol. 5, no. 1, pp. 30--42, 2023.

\bibitem{postgres_acid}
PostgreSQL Global Development Group, "PostgreSQL 16 Documentation," 2025. [Online]. Available: \url{https://www.postgresql.org/docs}. Accessed: Apr. 26, 2026.

\bibitem{react_ts}
React Team, "React Documentation," 2025. [Online]. Available: \url{https://react.dev}. Accessed: Apr. 26, 2026.

\bibitem{tailwind}
Tailwind Labs, "Tailwind CSS Documentation," 2025. [Online]. Available: \url{https://tailwindcss.com/docs}. Accessed: Apr. 26, 2026.

\bibitem{docker_docs}
Docker Inc., "Docker and Docker Compose Documentation," 2025. [Online]. Available: \url{https://docs.docker.com}. Accessed: Apr. 26, 2026.

\bibitem{k8s_docs}
The Kubernetes Authors, "Kubernetes Documentation," 2025. [Online]. Available: \url{https://kubernetes.io/docs}. Accessed: Apr. 26, 2026.

\bibitem{smtp_practice}
A. K. Jain and R. Menon, "Reliable OTP delivery pipelines for high-volume transactional systems," \emph{Intl. J. Secure Comput.}, vol. 14, no. 3, pp. 201--214, 2023.

\bibitem{p2p_risk}
L. Wang, M. Brown, and K. Shah, "Risk controls in peer-to-peer transfer systems," \emph{IEEE Trans. Dependable Secure Comput.}, vol. 20, no. 4, pp. 1519--1533, 2024.

\bibitem{ux_fintech}
D. Kim and E. Park, "Human-centered interface patterns for mobile fintech trust," \emph{Proc. CHI Fintech Workshop}, pp. 11--18, 2023.

\bibitem{ledger_consistency}
J. Alvarado and T. Green, "Consistency models for digital wallet ledgers in distributed services," \emph{ACM Trans. Internet Technol.}, vol. 24, no. 1, pp. 1--24, 2024.

\end{thebibliography}

\end{document}
